% titel: Blatt 3, A6: Verdünnen (Massen- und Stoffmengenkonzentration)
% Teil II: der durchgerechnete Weg.

\subsection*{Gegeben und gesucht}
Ziel: $V = \qty{50}{mL}$ einer Lösung mit
$c(\mathrm{HNO_3}) = \qty{6}{mmol/L}$. Stammlösung mit
$b = \qty{750}{mg/L}$. Gesucht: das mit der Pipette zu entnehmende
Volumen $V_0$.

$M(\mathrm{HNO_3}) \approx \qty{63.01}{g/mol}$.

\subsection*{Rechnung}
1. Ziel-Stoffmenge:
\begin{equation}
  n = c\,V = 6\times10^{-3}\ \frac{\mathrm{mol}}{\mathrm{L}}
  \times 0.050\ \mathrm{L} = 3.00\times10^{-4}\ \mathrm{mol}.
\end{equation}
2. Konzentration der Stammlösung in $\mathrm{mol/L}$ umrechnen:
\begin{equation}
  c_0 = \frac{b}{M}
  = \frac{0.750\ \mathrm{g/L}}{63.01\ \mathrm{g/mol}}
  = 1.190\times10^{-2}\ \mathrm{mol/L}.
\end{equation}
3. Entnahmevolumen:
\begin{equation}
  V_0 = \frac{n}{c_0} = \frac{3.00\times10^{-4}}{1.190\times10^{-2}}
  = 2.52\times10^{-2}\ \mathrm{L} = \qty{25.2}{mL}.
\end{equation}

\begin{ergebnis}
Es müssen $V_0 \approx \qty{25.2}{mL}$ der Stammlösung entnommen und auf
$\qty{50}{mL}$ aufgefüllt werden.
\end{ergebnis}

\begin{kontrolle}
Massen-Check: $3.00\times10^{-4}\ \mathrm{mol}$ entsprechen
$3.00\times10^{-4}\times 63.01 = 1.89\times10^{-2}\ \mathrm{g} =
\qty{18.9}{mg}$. Bei $b = 0.750\ \mathrm{g/L}$ liefern
$25.2\ \mathrm{mL}$ gerade $0.750\times 0.0252 = 1.89\times10^{-2}\ \mathrm{g}$
--- konsistent. Ausserdem ist $25.2\ \mathrm{mL} < 50\ \mathrm{mL}$,
also eine echte Verdünnung.
\end{kontrolle}

\begin{merksatz}
Beim Verdünnen bleibt $n = cV$ konstant; Konzentrationen zuerst auf
gemeinsame Einheiten bringen (über $c = b/M$), dann $V_0 = n/c_0$.
\end{merksatz}
